% =============================================================
%  cv_joseph_wallace.tex  —  v3
%  Font:  newtxtext (Times-based serif, clean academic look)
%  Build: latexmk -pdf cv_joseph_wallace.tex
% =============================================================

\documentclass[9pt, a4paper]{extarticle}

% ---- Encoding & font ----------------------------------------
\usepackage[T1]{fontenc}
\usepackage[utf8]{inputenc}
\usepackage{newtxtext}
\usepackage{microtype}

% ---- Layout -------------------------------------------------
\usepackage[top=0.9cm, bottom=0.9cm, left=1.8cm, right=1.8cm]{geometry}
\usepackage{enumitem}
\usepackage{tabularx}
\usepackage{hyperref}

\hypersetup{
  colorlinks = true,
  urlcolor   = black,
  linkcolor  = black,
  citecolor  = black,
  pdfauthor  = {Joseph Wallace},
  pdftitle   = {CV -- Joseph Wallace},
}

\pagestyle{empty}
\setlength{\parindent}{0pt}
\setlength{\parskip}{0pt}

% ---- Section heading: bold label + thin rule below ----------
\newcommand{\cvsection}[1]{%
  \par\vspace{6pt}%
  \noindent{\normalsize\bfseries #1}%
  \par\vspace{1pt}%
  \noindent\rule{\linewidth}{0.4pt}%
  \par\vspace{2pt}%
}

% ---- Entry: bold role + plain date right, plain org below ---
\newcommand{\cventry}[3]{%
  \par\vspace{2pt}%
  \noindent\begin{tabularx}{\linewidth}{@{}X r@{}}
    \textbf{#1} & #3 \\[-2pt]
    #2 & \\
  \end{tabularx}%
  \vspace{-3pt}%
}

% ---- Compact bullets ----------------------------------------
\newenvironment{cvbullets}{%
  \begin{itemize}[
    leftmargin = 1.4em,
    label      = \textbullet,
    topsep     = 1pt,
    itemsep    = 0pt,
    parsep     = 0pt
  ]%
}{%
  \end{itemize}\vspace{0pt}%
}

% =============================================================
\begin{document}
% =============================================================
%  HEADER  (left-aligned)
% =============================================================
\noindent{\fontsize{17}{21}\selectfont\bfseries Joseph Wallace}\\[2pt]
\noindent Computational Chemist\enspace|\enspace
           Molecular Modelling\enspace|\enspace
           Ligand Screening Workflows\\[3pt]
\noindent{\small
  \href{mailto:joseph.wallace@iit.it}{joseph.wallace@iit.it}%
  \enspace|\enspace
  \href{https://www.linkedin.com/in/joseph--wallace/}{LinkedIn}%
  \enspace|\enspace
  \href{https://github.com/joewa6}{GitHub}%
  \enspace|\enspace
  \href{https://joewa6.github.io/}{Web}%
  \enspace|\enspace
  Genoa, Italy%
}

\vspace{3pt}
\noindent\rule{\linewidth}{0.6pt}

% =============================================================
%  PROFESSIONAL SUMMARY
% =============================================================
\cvsection{Professional Summary}

Computational chemist with a background in molecular modelling, simulation,
and structure-focused analysis across biomolecular and nanoscale systems.
Currently building reproducible ligand-screening, docking, and rescoring
workflows for drug-discovery-style prioritisation, with a strong focus on
physically grounded modelling and benchmark-aware evaluation.

% =============================================================
%  CORE SKILLS
% =============================================================
\cvsection{Core Skills}

\noindent\textbf{Molecular Modelling \& Screening:} Docking workflows,
rescoring strategy, molecular interaction analysis, structure-based modelling,
comparative ranking and benchmark evaluation.

\noindent\textbf{Scientific Computing \& Data Analysis:} Python, NumPy,
pandas, SciPy, matplotlib, reproducible analysis pipelines, workflow
automation, data visualisation.

\noindent\textbf{Simulation Methods:} Molecular dynamics, enhanced sampling,
umbrella sampling, metadynamics, Hamiltonian replica exchange,
free-energy-oriented analysis.

\noindent\textbf{Cheminformatics \& Molecular Tools:} RDKit, PyMOL, VMD, xTB.

\noindent\textbf{HPC \& Research Computing:} Linux environments, HPC workflows,
job submission, simulation pipeline management. Software: GROMACS, PLUMED.

% =============================================================
%  PROFESSIONAL EXPERIENCE
% =============================================================
\cvsection{Professional Experience}

\cventry{Postdoctoral Researcher}%
        {University of Padova / Italian Institute of Technology, Italy}%
        {2022--Present}
\begin{cvbullets}
  \item Lead computational modelling and analysis on nanoparticle--molecule
        systems; designed Python-based pipelines for automated processing,
        visualisation, and benchmark-style comparison across systems.
  \item Applied molecular simulation and enhanced-sampling methods to study
        structure, dynamics, and interaction behaviour in complex molecular
        systems, translating results into interpretable and decision-ready
        outputs.
\end{cvbullets}

\cventry{PhD Research Attachment}%
        {Universit\'{e} Sorbonne Paris Nord, France}%
        {2021--2022}
\begin{cvbullets}
  \item Supported research continuity during laboratory relocation; contributed
        to computational research within an international academic environment.
\end{cvbullets}

\cventry{PhD Research Attachment}%
        {A*STAR Institute of High Performance Computing, Singapore}%
        {2019--2021}
\begin{cvbullets}
  \item Conducted molecular simulation and structural analysis research in an
        HPC environment, with emphasis on scalable and reproducible workflows.
\end{cvbullets}

% =============================================================
%  EDUCATION
% =============================================================
\cvsection{Education}

\cventry{PhD, Computational Biochemistry}%
        {University of Liverpool, UK}%
        {2017--2022}
\noindent\textit{Thesis:} How does constraining folded peptides within dense
self-assembled monolayers on gold nanoparticles affect their structure?
Combined molecular simulation and experimental data to study structure,
dynamics, and interactions in peptide--nanoparticle systems.

\cventry{BSc Biochemistry, First Class Honours}%
        {University of Liverpool, UK}%
        {2014--2017}
\noindent\textit{Project:} Solving the crystal structure of Cytochrome
bc\textsubscript{1} for the development of antimalarials.

% =============================================================
%  SELECTED PROJECTS
% =============================================================
\cvsection{Selected Projects}

\noindent\textbf{Docking + Rescoring Benchmark for Ligand Ranking}
\begin{cvbullets}
  \item Built a reproducible workflow to test whether rescoring improves
        prioritisation over baseline docking on a curated protein--ligand
        benchmark.
  \item Focused on ranking analysis, benchmark design, and explicit failure-mode
        evaluation rather than inflated performance claims.
  \item \textit{Links:}
        \href{[PROJECT_LINK_PLACEHOLDER]}{Project page}%
        \enspace|\enspace
        \href{[REPO_LINK_PLACEHOLDER]}{Repository}
\end{cvbullets}

\noindent\textbf{Simulation Analysis Pipelines for Molecular Systems}
\begin{cvbullets}
  \item Designed Python-based workflows for processing, comparing, and
        visualising molecular simulation outputs across multiple systems.
  \item Emphasis on reproducibility, interpretable analysis, and workflow reuse.
\end{cvbullets}

% =============================================================
%  SELECTED PUBLICATIONS
% =============================================================
\cvsection{Selected Publications}

\begin{enumerate}[leftmargin=1.6em, label=\arabic*.,
                  topsep=0pt, itemsep=3pt, parsep=0pt]

  \item Wallace, J., Riccardi, L., Mancin, F., \& De~Vivo, M. (2025).
        A Scoring Function for Monolayer-Protected Gold Nanoparticles Capable
        of Recognizing Small Organic Molecules in Solution.
        \textit{Journal of Chemical Theory and Computation}, \textbf{21}(21),
        11070--11078.
        \href{https://doi.org/10.1021/acs.jctc.5c01278}{DOI:~10.1021/acs.jctc.5c01278}

  \item Franco-Ulloa, S., Cesari, A., Zanoni, G., Riccardi, L., Wallace, J.,
        et al. (2025).
        Rational design of gold nanoparticle-based chemosensors for detection
        of the tumor marker 3-methoxytyramine.
        \textit{Chemical Science}, \textbf{16}, 6282--6289.
        \href{https://doi.org/10.1039/D4SC08758E}{DOI:~10.1039/D4SC08758E}

  \item Wallace, J. (2022).
        \textit{How Does Constraining Folded Peptides within Dense
        Self-Assembled Monolayers on Gold Nanoparticles Affect Their Structure?
        An Experimental and Computational Investigation.}
        PhD thesis, University of Liverpool.

\end{enumerate}

\end{document}
